\chapter{A Distributed Unified Architecture}
\label{ch:dua}
\doublespacing
\vspace{-2.5cm}

The aim of this Chapter is to present the core contribution of this work. The introductory Chapters defined the problem setting, outlining a situation in which distributed autonomous systems are becoming ubiquitous in many aspects of everyday life, but at the same time increasingly complex to develop and manage. From the researcher's and the developer's perspective, a versatile framework to work on such systems is needed, but once some specific yet broad requirements are stated one finds that the current state of the art does not provide suitable solutions. Still, if the trends from the last decade in software engineering and computer science are examined, one can find that many tools have been developed that address specific portions of the original problem.

Recalling once again the leading principle of this entire work, we are about to build a new type of wheel without starting completely from scratch: we have some good tools ready at hand, but most importantly some pieces are already available, whereas some others are still missing. Before we begin, however, it is important to further clarify the actual setting in which this work has been developed, in particular with respect to the Hardware (H) requirement from \Cref{sec:prb_requirements} and to supported operating systems: the scope of the design process has been to build a framework that can be adapted to many target platforms but naturally, due to availability and time constraints, the actual implementation at the time of writing has been tested on a limited yet representative set of target platforms. This process has been driven by the need to provide a proof of concept for the framework, and at the same time to demonstrate its feasibility and effectiveness in real-world scenarios over the years in which this work has been carried out. Nonetheless, the chosen set of target platforms is broad enough that the resulting framework works on all of them and can be adapted, with minimal effort and specific changes, to other platforms as well, basically repeating the same setup process that has been followed for the current set of platforms. Furthermore, as clarified in the following, other assumptions stand regarding the configuration of the host operating system: for simplicity, and to benefit from the existing solutions introduced previously, the framework has been designed to work on systems that support the solutions described in \Cref{sec:version_control,sec:containerization}. Having defined this starting point, the framework follows by integrating what is ready and then developing what is missing, to first provide a system abstraction layer as a solid foundation and then to build an architectural template that can be used to develop and deploy modules and software for distributed autonomous systems satisfying the original requirements. Additional tools can subsequently be developed and integrated to further enhance the framework, addressing issues that are specific to some use cases like, \emph{e.g.}, the ones described later on in this work.

The organization of this Chapter reflects this description, which is nothing but the recollection of the design and development process that led to the \emph{Distributed Unified Architecture} (DUA).

\section{Assumptions}
\label{sec:dua_assumptions}
Before delving into the details of the system abstraction layer offered by DUA, it is instrumental to describe the host system configuration that it expects. Recalling what has been said at the beginning of this Chapter, the following brief set of assumptions hold mainly for the sake of simplicity: satisfying them allows to employ solutions that are readily available, represent current industry standards in their respective fields, and help to solve significant portions of the problem at hand. What follows from them constitutes a working proof of concept, but nothing prevents the framework from being adapted to different configurations should any requirement change, or should the need arise to support different host system configurations. For now, what is provided in this work fullfills the initial requirements, and any expansion or adaptation is left as future work.

Moreover, for clarity, the following definition is in order here: the term \emph{host system configuration}, in this context, refers to a combination of hardware and software that identifies a specific class of systems; each class shall offer and support a defined set of functionalities, hardware components, and operating systems.

\begin{description}
  \item[Hardware platform type] The framework is designed to work on systems that range from single-board computers to more powerful general-purpose computers or even servers, of any architecture, provided that they support the solutions described in \Cref{sec:version_control,sec:containerization}. This assumption excludes lower-end systems like, \emph{e.g.}, microcontrollers, which do play important roles in the field of autonomous systems. This choice has been motivated by two facts: first, microcontrollers and similar systems usually have very limited and constrained sets of resources, which make them incompatible with the software solutions required to build and manage the proposed framework. Second, the variety of microcontroller architectures available today is so vast and each of them is so different from the others, also in terms of operating system support on platforms where this is available, that extending the proposed approach to this class of systems would require a significant effort. This is not to say that the proposed framework cannot be adapted to work with microcontrollers, but that this is left as future work, since their integration would certainly require a separate design process.
  \item[Support for version control] The host system must support a version control system, as described in \Cref{sec:version_control}. This is a fundamental requirement, as the framework relies on version control to manage the software modules that are developed and deployed on the components of the distributed system. The choice of Git as the version control system is motivated by its many features that make it highly efficient when working with projects that embed many modules, as described in \Cref{sec:git}.
  \item[Support for containerization] The host system must support containerization, as described in \Cref{sec:containerization}. This is another fundamental requirement, as the framework relies on containerization to isolate and manage software components. The choice of Docker as the containerization solution is motivated by its widespread adoption and the functionalities that it offers, enabling it to build replicable environments on a multitude of platforms, as described in \Cref{sec:docker}.
  \item[Operating system support] The framework is designed to natively work on the Linux operating system, as it is the most widely used operating system in the field of distributed autonomous systems. This choice is motivated by the fact that Linux, in its many distributions, offers the widest support for hardware devices and software tools. Moreover, it is built to be developed onto, being natively open-source and offering a stable and fully equipped environment right from the start for the development of software of any kind, especially that which is intended to directly access any hardware resource connected to the host machine. Furthermore, Linux is the operating system that best supports both version control and containerization solutions, whose importance in this project has been highlighted in the previous points. For these same reasons, when a SBC is released to the market, it is almost certain that its natively supported operating system is a Linux distribution, so to offer a consistent environment across all platforms supported by the framework it is natural to assume that the host configurations share this level of similarity. Among the many possible Linux distributions to choose from, this work does not impose any specific one but Ubuntu Linux in one of its latest Long-Term Support (LTS) versions available at the time of writing\footnote{22.04, 24.04.} has been employed for the development and testing of the framework, and is used as the reference for the system abstraction layer described later on. Again, this choice is motivated by similar reasons: Ubuntu is the most popular and easy-to-use Linux distribution, and it is widely supported by many software tools and packages as well as by many SBC vendors. As in any other such case, the use of a different distribution is possible provided that one manages to obtain a working environment with Git, Docker, and any other software dependency that they might require. To conclude, let us note that this assumption is not a strict requirement since the framework can be adapted to work also on Windows and macOS operating systems thanks to Docker, but it is recommended to use a Linux distribution since some features that are crucial for this work are not supported on those OSes due to their different support for containerization which, as explained in \Cref{sec:docker}, always involves a layer of virtualization. Such features will be introduced and described in the following sections, together with their role in the proposed framework but for now, since they mainly concern containerization and map to specific Docker functionalities, a summary table is provided in \Cref{tab:docker_features} in which each feature appears with the name of the corresponding Docker functionality. For completeness, note that while macOS offers a good support to containers running a Linux-based image thanks to Docker Desktop, the best way to run them on Windows is currently to use Docker Desktop in conjunction with the Windows Subsystem for Linux (WSL) \cite{microsoft_windows_nodate}.
  \begin{table}[h]
    \centering
    \begin{tabular}{|c|c|c|c|}
        \hline
        & \textbf{Linux} & \textbf{Windows} & \textbf{macOS} \\
        \hline
        \textbf{Volume mounts} & \checkmark & \checkmark & \checkmark \\
        \hline
        \textbf{Host networking} & \checkmark & \xmark & \xmark \\
        \hline
        \textbf{Graphical applications} & \checkmark & \checkmark & \xmark \\
        \hline
        \textbf{Hardware devices} & \checkmark & \xmark & \xmark \\
        \hline
        \textbf{Capabilities} & \checkmark & \xmark & \xmark \\
        \hline
        \textbf{IPC namespaces} & \checkmark & \xmark & \xmark \\
        \hline
    \end{tabular}
    \caption{Docker features relevant for the DUA framework and their support on different operating systems.}
    \label{tab:docker_features}
  \end{table}
\end{description}

\section{System abstraction layer: dua\_foundation}
\label{sec:dua_foundation}
The most important problem solved by DUA is also the first that it addresses: providing a similar environment across multiple system configurations. Achieving this is not a trivial task: the hardware platforms and operating systems that are of interest, and mentioned in \Cref{sec:dua_assumptions}, are very different from each other; yet, the aim is to provide an environment that offers the same OS APIs, libraries, and tools across all supported platforms. Nonetheless, DUA operates within the limits of what is possible without unnecessarily incresing the overhead: some aspects of the underlying hardware are not abstracted on purpose since it is not necessary, as shown later on, to provide a consistent experience, but instead would imply introducing significant overheads in the execution of any application. This is the case of, \emph{e.g.}, the hardware architecture: as illustrated in the following, currently supported platforms include solutions based on both the x86 and ARM architectures, and the framework is designed to work on both of them without hiding them or their differences under an abstraction layer. This is due to multiple factors: first, there is no practical necessity to do so, and second, the developer using the framework can benefit of the peculiarities and specific features of each platform on which they want to deploy their software. Last but not least, abstracting the hardware architecture would imply using hardware emulators such as QEMU \cite{bellard_qemu_nodate}, which are very powerful but also introduce significant overheads in the execution of any application, are not always available on any platform, and complicate access to system resources such as hardware devices which are instead of paramount importance in the practical contexts of this work, as shown in \Cref{tab:docker_features}.

This first part of DUA is called \texttt{dua\_foundation} \cite{masocco_dua-foundation_nodate}. The starting intuition is that environments defined by a set of OS APIs and software tools can be easily replicated and managed in the form of Docker images, as described in \Cref{sec:docker}. Building a Docker image tailored to every supported platform would allow to account for the specific characteristics of each one, but providing a consistent system configuration inside containers that use such images. Each researcher or developer, \emph{i.e.}, each \emph{user} of the DUA framework, would then have to add their own software dependencies and packages to work on their specific modules. Let us call a user-developed software module based on DUA, a \emph{unit}. To define the system configuration that a unit requires, the Docker image is still the best tool available for the same reasons. For the sake of consistency along the whole chain, Docker images related to units would have to start from the same set of base images as explained in \Cref{sec:docker}, which are the ones built by \texttt{dua\_foundation}; the specific one is chosen depending on the hardware platform to work on. Let us then refer to images built by \texttt{dua\_foundation} as to \emph{base units}.

The way a base unit is created differs even significantly for each, but the setup process they follow is the same for all. They all start from a platform-specific base image available on public registries like Docker Hub. Then, this image is enriched by a set of system packages which represent common dependencies in this field like, \emph{e.g.}, compilation toolchains. Finally, some more libraries are installed that are specific to the currently evaluated application areas of DUA: even if they may not be required by all units, they are included in the base units due to the frequency of their use in real-world scenarios, \emph{i.e.}, it would then be very likely to find different units that require them. At the end of the build process, each image is pushed on Docker Hub \cite{masocco_dotxautomationdua-foundation_nodate} to be available to all users of the framework.

The next Sections explain the contents of \texttt{dua\_foundation}: first, the set of software dependencies to be installed in every base unit is discussed, to describe what is the final system configuration achieved by every base unit. Then, currently supported hardware platforms are briefly illustrated, to show the variety of hardware that DUA integrates but also provide a first hint of the differences that the framework deals with. Finally, currently available base units are explained in more detail, together with the setup process that they follow.

\subsection{Common software dependencies}
\label{sec:dua_common_dependencies}
Providing a full list of common software dependencies that are provided in every base unit here would be a daunting task, as it would include a vast number of packages and libraries. The interested reader can directly inspect the Dockerfiles in \cite{masocco_dua-foundation_nodate}. However, to give a general idea of the environment that base units build, as well as to facilitate the inspection of their Dockerfiles, a summary of the most important software dependencies is provided here.

Each base unit starts from a set of system packages, that is, software packages that can be installed using the package manager of the Linux distribution that the base unit is built on, and provide basic functionalities. Since the distribution of choice for every base image, as anticipated, is Ubuntu Linux, the package manager is \texttt{apt}. The first set of installed packages includes a GCC compilation toolchain \cite{free_software_foundation_gcc_nodate}, a Python system environment, a Java development kit \cite{oracle_java_nodate}, shell utilities, and a set of libraries that are required in the subsequent steps and that may depend on the specific base unit, as explained later on. These are installed now, all in a single command, to reduce the number of layers in the Docker image and optimize its size, a well-established practice when building Docker images.

Next come the libraries that are specific to the application areas of DUA. Usually, the first one to be installed among those is the ROS 2 middleware, presented in \Cref{sec:ros2}. The amount of ROS 2 packages installed depends on the target platform of a base unit: if the target is a development PC, a wide variety of packages is installed including simulators and visualizers, whereas if the target is an SBC, a reduced set of packages is installed skipping, \emph{e.g.}, visualization tools; each ROS 2 installation, however, includes both the eProsima Fast DDS \cite{eprosima_fast_nodate} and the Eclipse Cyclone DDS \cite{eclipse_foundation_cyclone_nodate} implementations of the DDS middleware presented in \Cref{sec:dds}, as well as the Zenoh middleware illustrated in \Cref{sec:zenoh}. Then is the time for the OpenCV computer vision library \cite{opencv_team_opencv_nodate,culjak_brief_2012} and the Eigen linear algebra library \cite{jacob_eigen_nodate}. The former is currently the de-facto standard in open-source computer vision, originally developed by Intel in the late 90s and then expanded by a vast team of developers, it now offers a wide range of functionalities and the implementation of thousands of algorithms ranging from image processing to mathematical operators. The latter is a C++ template library for linear algebra, which is widely used in the field of robotics and computer vision for its efficiency and ease of use. Finally, the Gazebo simulator presented in \Cref{sec:ros2_sim} is installed.

The setup of a base unit concludes with specific tools that vary from time to time, depending on the evolution of application scenarios and supported platforms. Among these, the most important is currently the \texttt{dua\_utils} collection, which is part of DUA and is illustrated in {\color{red}TODO dua\_utils sec. ref.}. Their position at the end of the setup process is intended: these are the last stages of each Dockerfile contained in \cite{masocco_dua-foundation_nodate}, hence they correspond to the last layers in the image. This is done to allow to easily change these stages in case some tools are not needed or some new ones are to be added, without having to rebuild the whole image from scratch but only the last layers. This is another common practice in Docker image building, made possible by the ability of Docker to \emph{cache} image layers among subsequent builds, and it is used here to optimize the time required to build a base unit.

\subsection{Currently supported hardware platforms}
\label{sec:dua_supported_platforms}
To develop a framework that could be applicable to a wide variety of practical cases, a good variety of hardware platforms has been considered for the initial implementation. The goal of this Section is to present them in due detail, to show the variety of hardware that DUA integrates but also provide a first hint of the differences that the framework deals with. The categorization of the platforms made here is based on their system resources and intended use: it is not a strict classification, but a way to group them to better understand the differences among them. Note that, for simplicity, only 64-bit architectures are currently supported by DUA, as they have become the standard in the last decade and are now common in the microprocessor market.

\subsubsection{General-purpose PCs}
\label{sec:dua_supported_platforms_pc}
In the context of this work, this is the simplest category to address and also the most common. The first base units, as well as the other pieces of the DUA framework, were put together and tested on this kind of platforms to ensure that the main concepts could work. By \emph{general-purpose PC}, we mean a computer that is not specifically designed for a particular task, but that can be used for a wide variety of applications. This category includes desktop computers, laptops, and servers. In this context, it is assumed that among this category of systems, the ones that are of interest for this work are based on either the x86-64 or the ARM64 hardware architectures. In the practical and applicative contexts relevant to the DUA framework, these systems are used for development and testing purposes, as well as for running simulations and visualizations, or intensive real-time computations. The software tools that are used on these systems are the same as the ones that are used on the other platforms, but the hardware resources available on them are usually more powerful: many CPU cores, RAM, and storage space are available. Dedicated coprocessors like, \emph{e.g.}, GPUs may also be available on these platforms.

\subsubsection{Single-board computers}
\label{sec:dua_supported_platforms_sbc}
In the field of autonomous systems, and especially in robotics, SBCs are very common. They are used in many applications, from small robots to drones, and they are also used in many research projects. The main reason for their popularity is that they are very cheap, small, and easy to use and integrate in larger systems, also thanks to their multiple communication interfaces. They are also very versatile, as they can be used for many different tasks. In the context of this work, SBCs are intended to be used to run the software modules that make an autonomous system perform its tasks. The hardware resources available here are generally more limited than the previous case, and GPUs and similar dedicated solutions may also not be available. This is also what makes SBCs so versatile in this field: their variety allows the system designer to choose the best solution that fits all the requirements, from software support, to mechanical constraints and energy consumption. In this work and for the explanatory purposes of this Section, three classes of single-board computers are considered, making broad distinctions based on their hardware architecture and intended use.

The first class consists of single-board computers based on the ARM64 architecture. They are probably the most common and versatile type of SBC, due to some reknown features of the ARM architectures. First, ARM architectures usually have a very low power consumption, which makes them ideal for battery-powered applications. Second, their open design allows for a wide variety of implementations, which makes them very versatile and has paved the way for some very small but powerful implementations, perfect for SBCs and embedded systems. Last, ARM architectures have a very tidy instruction set, which contributes to both their low power consumption and high efficiency.
